% AI-written, AI cross-reviewed; no human sign-off; not a claim that RH is proved
% Atlas unit T18. Local manuscript v3, item 5.28, heading on printed PDF p. 45.
% PositivityTheory.Pf1.HilbertSamuelMoriwaki.hsUse_outright
% paper2/split_p012.tex:317–335. Verbatim excerpt; SHA-256 2ddd309ab91adcc1615c47a8ffc382ab074a7400993be5524602b07463623927.
% Original TeX uses the macros and cross-reference keys of the source manuscript.


% Start of atlas edition annotation. The following block is outside the original excerpt.
\paragraph{Status of named inputs in this reading.}
\noindent\nolinkurl{PositivityTheory.ResPlug.FromRes.resK} --- in-repo constant (\path{Development/WeilBound/CoreStatements/ResolutionCharZero/PlugResK.lean:7--8}); not compared.\par
\noindent\nolinkurl{PositivityTheory.Tls2.Pl.poincareLelong} --- in-repo constant (\path{Development/WeilBound/CoreStatements/TemperedRepresentationTheory/PoincareLelong/PoincareLelongInv.lean:13--17}); not compared.\par
\noindent\nolinkurl{PositivityTheory.MbAis.DenseCurves.main} --- in-repo constant (\path{Development/WeilBound/CoreStatements/AutomorphicKernels/MoriwakiBigness/MoriwakiSymmetry/EqMapSymm.lean:23--119}); not compared.\par
\noindent\nolinkurl{GlobalTheory.T.arith_intersection_system_of_poincare_lelong} --- in-repo constant (\path{Development/AutomorphicSpectrum/Proofs/MetricProofs/ArithIntersectionSystemOfPoincareLelongStatement.lean:11--12}); not compared.\par
% End of atlas edition annotation.

\begin{formalstatement}[Three arithmetic Hilbert--Samuel consequences]\label{sub:T18}

For each intersection system and finite multihomogeneous presentation as
above the following three statements hold.
(i) On an arithmetic variety of dimension $D+1$, relative semipositivity,
fibre nefness, conjugation invariance and nonnegative degree on every
arithmetic curve imply $I_P(A^{D+1})\ge0$.
(ii) For a birational characteristic-zero presentation with $D\ge1$,
small-generation data and conjugation invariance, the limsup of the
small-section count for \nolinkurl{chartLat}, divided in degree $n$ by
$n^{D+1}/(D+1)!$, is at most $I_P(A^{D+1})$.
(iii) For a birational presentation over a field with $D\ge1$, nonnegative
multidegree, relative semipositivity at $P_0\subseteq P$, bounded monomials
and conjugation invariance, the same limsup bound holds for
\nolinkurl{chartLatNm}.

\end{formalstatement}\subnode{T18}

\noindent Corollary 5.5 gives arithmetic Hilbert--Samuel for a nef line. The three formal clauses involve two different chart lattices and a curve test. \href{https://arxiv.org/pdf/math/0612269}{A. Moriwaki, Continuity of volumes on arithmetic varieties, Corollary 5.5}~\cite{r21-moriwaki}.
